% ============================================================
\chapter{SQL --- Aggregation and Window Functions}
\label{ch:aggregation}
% ============================================================

Until now, our SQL queries returned individual rows. But often, we want summaries: the total number of students, the average grade by department, the revenue per quarter. This is the role of \emph{aggregate functions} (\texttt{COUNT}, \texttt{SUM}, \texttt{AVG}, \texttt{MAX}, \texttt{MIN}) combined with \texttt{GROUP BY}. \emph{Window functions} (\texttt{OVER}), introduced in SQL:2003, go further: they compute aggregates while preserving individual rows---rank, cumulative sum, moving average---without requiring complex subqueries.

\section{Aggregate Functions}

\begin{definition}[Aggregate Function]
An \emph{aggregate function} operates on a set of rows and returns a single
value summarizing that set.
\end{definition}

\begin{keyformulas}[title=\textbf{Standard Aggregate Functions}]
\begin{center}
\begin{tabular}{lll}
\toprule
\textbf{Function} & \textbf{Description} & \textbf{Ignores NULL?} \\
\midrule
\texttt{COUNT(*)} & Number of rows & No \\
\texttt{COUNT(col)} & Number of non-NULL values & Yes \\
\texttt{COUNT(DISTINCT col)} & Number of distinct values & Yes \\
\texttt{SUM(col)} & Sum & Yes \\
\texttt{AVG(col)} & Average & Yes \\
\texttt{MIN(col)} & Minimum & Yes \\
\texttt{MAX(col)} & Maximum & Yes \\
\bottomrule
\end{tabular}
\end{center}
\end{keyformulas}

\begin{codeblock}[title=\textbf{Simple Aggregations}]
\begin{minted}{sql}
SELECT COUNT(*) AS num_enrollments,
       COUNT(DISTINCT student_id) AS num_students,
       AVG(grade) AS avg_grade,
       MIN(grade) AS min_grade,
       MAX(grade) AS max_grade,
       SUM(grade) AS total_grades
FROM Enrollments;
\end{minted}
\end{codeblock}

\begin{codeoutput}
\begin{verbatim}
num_enrollments | num_students | avg_grade | min_grade | max_grade | total_grades
----------------+--------------+-----------+-----------+-----------+-------------
8               | 5            | 83.4375   | 55.0      | 97.0      | 667.5
\end{verbatim}
\end{codeoutput}

\begin{warning}[title=\textbf{AVG and NULL Values}]
\texttt{AVG} ignores \texttt{NULL} values. The average of $\{70, \texttt{NULL}, 90\}$
is $80$, not $53.3$. If you want to treat NULLs as $0$, use
\texttt{AVG(COALESCE(grade, 0))}.
\end{warning}

\section{GROUP BY --- Grouping}

\begin{definition}[GROUP BY]
The \texttt{GROUP BY} clause partitions rows into groups by one or more columns.
Aggregate functions are then applied to each group separately.
\end{definition}

\begin{keyformulas}[title=\textbf{Execution Order with GROUP BY}]
\begin{enumerate}
    \item \texttt{FROM} / \texttt{JOIN}
    \item \texttt{WHERE} (filter before grouping)
    \item \texttt{GROUP BY}
    \item \texttt{HAVING} (filter after grouping)
    \item \texttt{SELECT}
    \item \texttt{DISTINCT}
    \item \texttt{ORDER BY}
    \item \texttt{LIMIT}
\end{enumerate}
\end{keyformulas}

\begin{codeblock}[title=\textbf{Average per Student}]
\begin{minted}{sql}
SELECT s.last_name, s.first_name,
       COUNT(*) AS num_courses,
       ROUND(AVG(e.grade), 2) AS avg_grade,
       MIN(e.grade) AS min_grade,
       MAX(e.grade) AS max_grade
FROM Students s
JOIN Enrollments e ON s.student_id = e.student_id
GROUP BY s.student_id, s.last_name, s.first_name
ORDER BY avg_grade DESC;
\end{minted}
\end{codeblock}

\begin{codeoutput}
\begin{verbatim}
last_name | first_name | num_courses | avg_grade | min_grade | max_grade
----------+------------+-------------+-----------+-----------+----------
Taylor    | Emma       | 1           | 97.0      | 97.0      | 97.0
Miller    | Claire     | 2           | 88.75     | 82.0      | 95.5
Brown     | Alice      | 2           | 82.25     | 76.0      | 88.5
Davis     | Bob        | 2           | 81.5      | 72.0      | 91.0
Wilson    | David      | 1           | 55.0      | 55.0      | 55.0
\end{verbatim}
\end{codeoutput}

\begin{antipattern}[title=\textbf{Column Outside GROUP BY Without Aggregation}]
Every column in the \texttt{SELECT} must either appear in \texttt{GROUP BY}
or be inside an aggregate function:
\begin{minted}{sql}
-- ERROR: first_name not in GROUP BY or aggregated
SELECT last_name, first_name, AVG(grade) FROM Enrollments
JOIN Students ON ... GROUP BY last_name;

-- CORRECT:
SELECT last_name, first_name, AVG(grade) FROM Enrollments
JOIN Students ON ... GROUP BY last_name, first_name;
\end{minted}
MySQL in non-strict mode silently accepts this error (indeterminate result).
PostgreSQL and SQLite correctly reject it.
\end{antipattern}

\section{HAVING --- Filtering After Grouping}

\begin{definition}[HAVING]
The \texttt{HAVING} clause filters groups after aggregation. It is to
\texttt{GROUP BY} what \texttt{WHERE} is to \texttt{FROM}:
\texttt{WHERE} filters rows, \texttt{HAVING} filters groups.
\end{definition}

\begin{codeblock}[title=\textbf{Filtering with HAVING}]
\begin{minted}{sql}
-- Students with average above 80
SELECT s.last_name, AVG(e.grade) AS avg_grade
FROM Students s
JOIN Enrollments e ON s.student_id = e.student_id
GROUP BY s.student_id, s.last_name
HAVING AVG(e.grade) > 80
ORDER BY avg_grade DESC;

-- Courses with at least 2 enrolled students
SELECT c.title, COUNT(*) AS num_enrolled
FROM Courses c
JOIN Enrollments e ON c.course_id = e.course_id
GROUP BY c.course_id, c.title
HAVING COUNT(*) >= 2;
\end{minted}
\end{codeblock}

\begin{codeoutput}
\begin{verbatim}
-- Average > 80:
last_name | avg_grade
----------+----------
Taylor    | 97.0
Miller    | 88.75
Brown     | 82.25
Davis     | 81.5

-- Courses with >= 2 enrolled:
title      | num_enrolled
-----------+-------------
Databases  | 3
Algorithms | 2
Statistics | 2
\end{verbatim}
\end{codeoutput}

\section{Advanced Grouping}

\subsection{GROUPING SETS, ROLLUP, CUBE}

\begin{codeblock}[title=\textbf{ROLLUP --- Hierarchical Subtotals (PostgreSQL)}]
\begin{minted}{sql}
SELECT s.major,
       c.title,
       AVG(e.grade) AS avg_grade,
       COUNT(*) AS num
FROM Students s
JOIN Enrollments e ON s.student_id = e.student_id
JOIN Courses c ON e.course_id = c.course_id
GROUP BY ROLLUP (s.major, c.title)
ORDER BY s.major, c.title;
\end{minted}
\end{codeblock}

\texttt{ROLLUP} generates hierarchical subtotals: by (major, course),
by major, and the grand total.

\texttt{CUBE} generates all possible combinations of subtotals.

\section{Window Functions}

\begin{definition}[Window Function]
A \emph{window function} performs a calculation across a set of rows related
to the current row, \emph{without reducing the number of rows} in the result.
It uses the \texttt{OVER()} clause.
\end{definition}

\begin{theorem}[GROUP BY vs.\ Window Function]
\begin{itemize}
    \item \texttt{GROUP BY} + aggregation: reduces rows (one row per group).
    \item Window function: keeps all rows while adding a computed value.
\end{itemize}
\end{theorem}

\subsection{OVER Syntax}

\begin{keyformulas}[title=\textbf{Window Function Syntax}]
\begin{minted}{sql}
function(args) OVER (
    [PARTITION BY columns]
    [ORDER BY columns [ASC|DESC]]
    [ROWS BETWEEN start AND end]
)
\end{minted}
\end{keyformulas}

\subsection{Windowed Aggregation}

\begin{codeblock}[title=\textbf{Windowed Average by Major}]
\begin{minted}{sql}
SELECT s.last_name, s.major, e.grade,
       AVG(e.grade) OVER (PARTITION BY s.major) AS major_avg,
       e.grade - AVG(e.grade) OVER (PARTITION BY s.major) AS deviation
FROM Students s
JOIN Enrollments e ON s.student_id = e.student_id
ORDER BY s.major, s.last_name;
\end{minted}
\end{codeblock}

\begin{codeoutput}
\begin{verbatim}
last_name | major | grade | major_avg | deviation
----------+-------+-------+-----------+----------
Brown     | CS    | 88.5  | 76.5      | 12.0
Brown     | CS    | 76.0  | 76.5      | -0.5
Davis     | CS    | 72.0  | 76.5      | -4.5
Davis     | CS    | 91.0  | 76.5      | 14.5
Wilson    | CS    | 55.0  | 76.5      | -21.5
Miller    | Math  | 95.5  | 91.5      | 4.0
Miller    | Math  | 82.0  | 91.5      | -9.5
Taylor    | Math  | 97.0  | 91.5      | 5.5
\end{verbatim}
\end{codeoutput}

\subsection{ROW\_NUMBER, RANK, DENSE\_RANK}

\begin{definition}[Ranking Functions]
\begin{itemize}
    \item \texttt{ROW\_NUMBER()} --- unique sequential number within the partition.
    \item \texttt{RANK()} --- rank with gaps (ties receive the same rank; the
          next rank is incremented by the number of ties).
    \item \texttt{DENSE\_RANK()} --- rank without gaps.
\end{itemize}
\end{definition}

\begin{codeblock}[title=\textbf{Ranking Students by Grade}]
\begin{minted}{sql}
SELECT s.last_name, e.grade,
       ROW_NUMBER() OVER (ORDER BY e.grade DESC) AS row_num,
       RANK()       OVER (ORDER BY e.grade DESC) AS rnk,
       DENSE_RANK() OVER (ORDER BY e.grade DESC) AS dense_rnk
FROM Students s
JOIN Enrollments e ON s.student_id = e.student_id;
\end{minted}
\end{codeblock}

\begin{codeoutput}
\begin{verbatim}
last_name | grade | row_num | rnk | dense_rnk
----------+-------+---------+-----+----------
Taylor    | 97.0  | 1       | 1   | 1
Miller    | 95.5  | 2       | 2   | 2
Davis     | 91.0  | 3       | 3   | 3
Brown     | 88.5  | 4       | 4   | 4
Miller    | 82.0  | 5       | 5   | 5
Brown     | 76.0  | 6       | 6   | 6
Davis     | 72.0  | 7       | 7   | 7
Wilson    | 55.0  | 8       | 8   | 8
\end{verbatim}
\end{codeoutput}

\begin{codeblock}[title=\textbf{Top Student per Major (Top-N per Group)}]
\begin{minted}{sql}
WITH Ranked AS (
    SELECT s.last_name, s.major, e.grade,
           ROW_NUMBER() OVER (
               PARTITION BY s.major
               ORDER BY e.grade DESC
           ) AS rn
    FROM Students s
    JOIN Enrollments e ON s.student_id = e.student_id
)
SELECT last_name, major, grade
FROM Ranked
WHERE rn = 1;
\end{minted}
\end{codeblock}

\subsection{LAG, LEAD}

\begin{definition}[LAG and LEAD]
\begin{itemize}
    \item \texttt{LAG(col, n, default)} --- value of \texttt{col} $n$ rows
          \emph{before} the current row.
    \item \texttt{LEAD(col, n, default)} --- value of \texttt{col} $n$ rows
          \emph{after} the current row.
\end{itemize}
\end{definition}

\begin{codeblock}[title=\textbf{Comparing with the Previous Row}]
\begin{minted}{sql}
SELECT s.last_name, e.grade,
       LAG(e.grade, 1) OVER (ORDER BY e.grade DESC) AS prev_grade,
       e.grade - LAG(e.grade, 1) OVER (ORDER BY e.grade DESC) AS diff
FROM Students s
JOIN Enrollments e ON s.student_id = e.student_id
ORDER BY e.grade DESC;
\end{minted}
\end{codeblock}

\subsection{Sliding Windows (Frame)}

\begin{codeblock}[title=\textbf{Moving Average over 3 Rows}]
\begin{minted}{sql}
SELECT s.last_name, e.grade,
       AVG(e.grade) OVER (
           ORDER BY e.grade
           ROWS BETWEEN 1 PRECEDING AND 1 FOLLOWING
       ) AS moving_avg
FROM Students s
JOIN Enrollments e ON s.student_id = e.student_id
ORDER BY e.grade;
\end{minted}
\end{codeblock}

\begin{keyformulas}[title=\textbf{Window Frame Specifications}]
\begin{center}
\begin{tabular}{ll}
\toprule
\textbf{Specification} & \textbf{Description} \\
\midrule
\texttt{ROWS BETWEEN UNBOUNDED PRECEDING AND CURRENT ROW} & From start to current row \\
\texttt{ROWS BETWEEN 1 PRECEDING AND 1 FOLLOWING} & Sliding window of 3 rows \\
\texttt{ROWS BETWEEN CURRENT ROW AND UNBOUNDED FOLLOWING} & From current row to end \\
\texttt{RANGE BETWEEN ...} & Value-based (not position-based) \\
\bottomrule
\end{tabular}
\end{center}
\end{keyformulas}

\subsection{NTILE and Cumulative Percentages}

\begin{codeblock}[title=\textbf{Quartiles and Cumulative Percentage}]
\begin{minted}{sql}
SELECT s.last_name, e.grade,
       NTILE(4) OVER (ORDER BY e.grade DESC) AS quartile,
       ROUND(100.0 * SUM(e.grade) OVER (ORDER BY e.grade DESC
           ROWS BETWEEN UNBOUNDED PRECEDING AND CURRENT ROW)
           / SUM(e.grade) OVER (), 1) AS cumul_pct
FROM Students s
JOIN Enrollments e ON s.student_id = e.student_id
ORDER BY e.grade DESC;
\end{minted}
\end{codeblock}

\section{Python Integration}

\begin{codeblock}[title=\textbf{Aggregation and Window Functions with Python}]
\begin{minted}{python}
import sqlite3

conn = sqlite3.connect('university.db')
cur = conn.cursor()

# Aggregation with GROUP BY
cur.execute('''
    SELECT s.last_name, COUNT(*) as num, ROUND(AVG(e.grade), 2) as avg
    FROM Students s
    JOIN Enrollments e ON s.student_id = e.student_id
    GROUP BY s.student_id, s.last_name
    HAVING AVG(e.grade) > 80
    ORDER BY avg DESC
''')
print("Averages > 80:")
for row in cur.fetchall():
    print(f"  {row[0]:10} | {row[1]} courses | avg {row[2]}")

# Window function (SQLite >= 3.25)
cur.execute('''
    SELECT s.last_name, e.grade,
           RANK() OVER (ORDER BY e.grade DESC) as rnk
    FROM Students s
    JOIN Enrollments e ON s.student_id = e.student_id
''')
print("\nRankings:")
for name, grade, rnk in cur.fetchall():
    print(f"  #{rnk} {name:10} : {grade}")

conn.close()
\end{minted}
\end{codeblock}

\section*{Exercises}

\begin{exercise}
For each course, display the number of enrolled students, the average grade,
the minimum, and the maximum. Only show courses with at least 2 students.
\end{exercise}

\begin{exercise}
Rank students by their overall average and display their rank.
Use \texttt{DENSE\_RANK} to handle ties.
\end{exercise}

\begin{exercise}
For each enrollment, display the student's grade, their major's average,
and the deviation from that average. Use a window function.
\end{exercise}

\begin{exercise}
Display the top 2 grades per course. Use \texttt{ROW\_NUMBER} with
\texttt{PARTITION BY}.
\end{exercise}

\begin{exercise}
Calculate the cumulative sum of grades per student, ordered by course ID,
using a sliding window.
\end{exercise}
